\chapter{第五次“围剿”：南京政府全力以赴}

\section{国民党军的作战准备和作战方针}


从1930年代初开始，国共间连续展开数次“围剿”与反“围剿”战争，蒋介石和国民党军屡战不利，这其中，有着国内外多重方面的原因，但失利的阴影使蒋介石和国民党政权对中共不得不予正视。失败的痛苦使蒋痛切意识到：“我们纪律和工作超过土匪，超过共产党，然后才有剿清他们的希望……人家死中求生，拼命牺牲努力！我们苟且偷安，得过且过，似乎是要从生中避死，实在是死中不知求生！”\footnote{蒋介石：《剿匪应从精神、组织与纪律来奋斗》，《先总统蒋公思想言论总集》第11卷，台北，国民党中央党史委员会，1984，第20页。}在他耳提面命下，国民党内部的自我反省和对中共的研究明显加强，有关的研究、反思文献不断出笼。他们强调：“事实告诉我们，认定现在江西的土匪，与历史上的土匪是一样的容易消灭，这未免忽略了他的国际性和进步的伎俩。要是以为他怎样了不得，却又未免自减锐气。”\footnote{《由西成桥之役到下坪之役》，《汗血月刊》第1卷第5期，1934年7月20日。}其对中共的性质、力量和发展有了相对较为平心静气的估计。当时中共方面在接触到国民党这些研究成果后，敏锐感觉到“敌军屡受挫败之后，于战略战术上得了不少的进步”，断言国民党军此后的“‘剿共’计划，更见进步，更见周密，是可以预料的”\footnote{冰澈：《对于白军“剿共”的研究》，1931年编印，第3页。}。


国民党军的这种反思、检讨在1933年开始的庐山训练中有着鲜明的体现。第四次“围剿”失败后，蒋介石下决心对国民党军队进行全面整编，计划先精简编制，然后“轮流抽调三分之一的官长到后方，加以特别训练。大概每个官长训练半月至一月就行，在六七八这三个月之内，全部官长训练好了，等到九月，我们就可以将这种新编练的十六、七个师，照正式计划，来向匪区进剿”。\footnote{蒋介石：《剿匪基本工作之研究》，《先总统蒋公思想言论总集》第11卷，第42页。}1933年7月1日，南昌行营发出电令，指出：

\begin{quote}

土匪盘踞赣南，日形猖獗，迭经痛剿，未着特效。推厥原因，实由各部队中初级军官之武德、武学，尚欠深造所致。亟应严格训练，注入剿匪特要之学术科，以增进剿匪之效能。兹在庐山设立北路剿匪军官训练团，分期召集各部队中之中初级军官赴该团训练。\footnote{《庐山训练实纪》，赣粤闽湘鄂北路剿匪军军官训练团编印，1933，第58页。}
\end{quote}


根据这一思路，1933年7月军官训练团在江西庐山开办，到9月先后办了三期，受训者达7500余人。训练“惟一的目的，就是要消灭赤匪，所以一切的设施，皆要以赤匪为对象”。\footnote{蒋介石：《军官训练团训练的要旨与方法》，《庐山训练集》，南京新中国出版社，1947，第18页。}


对应着蒋介石关于“从前剿匪剿不了，并不是我们武力不够，而是我们精神不良”的认识，训练团特别注重战斗、团结精神的教育，耳提面命，树立各级长官的战斗意志和团结精神，以拉近和中共在这方面的距离。训练的结果，按蒋自己的说法：“因为时间过于短促，对于学术科没有多大的进步，但是各人的精神思想，和所表现的仪容、态度、动作，比两星期以前，完全不同了。”\footnote{蒋介石：《时间为一切事业与生命之母》，《庐山训练集》，第71页。}经过整训，国民党军队面貌确有所改观：“十一师自受创后到抚州补充训练，前师长萧乾以复仇雪耻为号召，提倡刻苦作风，到各连队与士兵同生活起居。为了练习长途行军，萧以身作则，脚穿草鞋，脚底抹了桐油，以作标榜。六十七师师长傅仲芳的行李只有半担，经常穿旧军衣，军中有伙夫头之称。霍揆彰、李树森等行军时都背米袋。”\footnote{杨伯涛：《蒋军对中央苏区第五次围攻纪要》，《文史资料选辑》第45辑，全国政协文史委员会编印，1964，第183页。}国民党军的这种改变，作为其对手的中共方面感受最深，第五次反“围剿”开始近半年后，经过往复交手，周恩来深有体会地写道：

\begin{quote}

蒋介石对于这些军官的训练，不能说是没有相当的结果，如果我们看到四次战争中白军军官的无能，那我们看到现在是狡猾机警得多了。他们懂得了如何防备我们打埋伏，如何避免运动战中整师整旅的被消灭，如何加强其侦查搜索与通信联络的工作，如何依靠堡垒与我们作战而很快的缩回堡垒去，这些都要算是他的进步。\footnote{周恩来：《五次战役中我们的胜利（一）——论持久战》，《红星》第33期，1934年3月18日。}
\end{quote}


同时，训练时干部集中、朝夕相处，这对来自各地方、各派别，曾经多次兄弟阋墙的军官“把眼前畛域派别的观念，和频年交相火并的宿怨前隙，不期然而然的消弭泯灭”，\footnote{《黄绍竑回忆录》，广西人民出版社，1991，第332页。}也发挥了一定作用。


以红军作为假想敌，郑重其事地开办庐山训练，反映蒋介石对再一次与红军作战的充分重视，如他所说：“此次剿匪，实关党国与本军之存亡，不可以大意轻易出之。”\footnote{《蒋介石1933年9月10日致熊式辉等电》，秦孝仪主编《中华民国重要史料初编——对日抗战时期》绪编（二），第397页。}这和他长沙被打下后仍称“长沙虽失，共犯实癣疥之疾耳”，\footnote{《蒋介石日记》，1930年8月5日。}以及第一、二次“围剿”时视红军为“地方事件”，\footnote{李家白：《反共第一次“围剿”的源头之役》，《文史资料选辑》第45辑，第76页。}仅出动一些杂牌部队有重大差异。吸取前几次“围剿”失败的教训，蒋介石将新一次“围剿”定位为军事、政治、经济、社会的总体战。1933年6月，蒋在南昌主持召开“剿共”军事会议，通过军事、谍报、宣传等“剿共”方案，确定第五次“围剿”的基本原则仍为前一年提出的“三分军事，七分政治”，即“用三分的力量作战，用七分的力量来推行作战区的政治”，\footnote{蒋介石：《剿匪成败与国家存亡》，张其昀主编《先总统蒋公全集》第1册，台北，中国文化大学，1984，第209页。}“一方面要发挥军事的力量，来摧毁土匪的武力；一方面要加倍地运用种种方法，消极地来摧毁土匪所有的组织，及在民众中一切潜势力”；“尤其是要教化一般民众，使他能倾向我们的主义，以巩固我们在民众中精神的堡垒”。\footnote{蒋介石：《推进剿匪区域政治工作的要点》，《先总统蒋公思想言论总集》第11卷，第234页。}


秉着总体战的思路，国民党方面采取了一系列政治、经济、社会政策，严密统治、收揽人心，其中，对苏区最具威胁的是封锁政策。通过修筑碉堡和大量设岗设哨，实行严密的经济、交通和邮电封锁，严禁粮秣、食盐、工业品和原材料等物资流入苏区，断绝其与外界的联系，蒋介石判断：“我们不仅在土地方面，经济方面，武器方面，军队方面的力量，统统超过他们几倍，乃至几十倍几百倍，就是专就壮丁的人数讲，可以说我们的补充无穷，而他们除了现存的这五六万人以外简直就再也没有了！”\footnote{蒋介石：《以自强的精神剿必亡的赤匪》，《先总统蒋公思想言论总集》第11卷，第127页。}“匪区数年以来，农村受长期之扰乱，人民无喘息之余地，实已十室九空，倘再予以严密封锁，使其交通物质，两相断绝，则内无生产，外无接济，既不得活动，又不能鼠窜，困守一隅，束手待毙。”\footnote{国民政府军事委员会委员长南昌行营：《处理剿匪省份政治工作报告》（下），南昌行营，1934，第十一章，第1～2页。}用心既狠且辣。


与对中共力量认识的变化相应，蒋介石的战略方针也一改前几次“围剿”常犯的急躁冒进错误，强调稳扎稳打，层层推进。总的作战指导方针是“不先找匪之主力，应以占领匪必争之要地为目的”，\footnote{蒋介石1933年10月17日战字二一三号训令，《赣粤闽湘鄂北路剿匪军第三路军五次进剿战史》（上），第三路军总指挥部参谋处，1937，第三章，第12页。}即以严密的工事和碉堡层层推进，通过缓慢但有效的占领方式，压缩红军作战区域，限制红军活动空间，迫使红军进行真面目的主力决战。这一战术的核心就是逼迫红军进行打资源、拼消耗、比人力的持久战，而在这三项因素中，国民党都占有着绝对优势。国民党军主攻部队第三路军作战方针明确规定：“本路军以消耗战之目的，采断绝赤区脉络、限制匪之流窜、打破其游击战术、封锁围进之策略……本战术取守势战略取攻势之原则，步步为营，处处筑碉，匪来我守，匪去我进。”\footnote{《赣粤闽湘鄂北路剿匪军第三路军作战计划》，《赣粤闽湘鄂北路剿匪军第三路军五次进剿战史》（上），第三章，第13页。}蒋介石在对参加“围剿”部队讲演时再三强调：“无论行军宿营或作战，也无论在匪区或非匪区，一定要随时随地，格外谨慎小心，严密防范，搜索，警戒，侦探，连络，和做防御工事，这几件事情，不可一刻懈怠。”\footnote{蒋介石：《剿匪必具的精神与当务之急》，《先总统蒋公思想言论总集》第11卷，第255页。}


蒋介石在第五次“围剿”中采取稳健持久的作战方针，实际也是出于对红军灵活机动战术及坚强战斗精神的惧怕和无奈，如他所说，红军“总是以迂为直，攻我不备，避实击虚，蹈瑕抵隙，他只要晓得我们哪一点力量单弱，哪一处防备不周，不管有几远距离，也不管是如何艰险的道路，就是集结主力来突破一点，或先派小的部队钻进来到处扰乱我们，以牵制我们主力作战，这就是他们的智谋，而‘超巅越绝，不畏险阻，耐饥忍渴，不避艰难’就是土匪唯一的惯技”。\footnote{蒋介石：《剿匪最重要的技能是什么》，《庐山训练集》，第124页。}在以快对快、以巧对巧的能力对抗中无法占据优势时，退而避敌锋芒，利用自身资源和人力优势一步步压迫对手，通过“逐渐消灭匪军的实力与资财”\footnote{《文件汇编》第三编，“赣粤闽湘鄂北路剿匪军第三路军”，1935，第1852页。}达到战胜的目的，是其无奈却不失明智的选择。为此，蒋明确告诉部下：

\begin{quote}

土匪用种种巧妙的方法来骚扰我们，我们只有先集结兵力，坚筑工事，用一个呆笨方法自己固守起来！以后再找好的机会来消灭他！土匪用声东击西，摇旗呐喊，以及种种虚声恫吓装模作样的巧妙方法来骚扰我们耳目，打击我们的精神，从而讨取便宜，我们只有一概不管，自己实实在在准备自己的实力，强固自己的工事，拿我们一切实在的东西，来对付土匪一切虚伪的花样，这就是所谓“以拙制巧，以实击虚”。\footnote{蒋介石：《主动的精义与方法》，《庐山训练集》，第196页。}
\end{quote}


在给陈诚的信中，他进一步谈道：

\begin{quote}

以匪之无远火器又无飞机高武器，不能妨碍我密集部队之行动，更可以匪之战术攻匪。故我军以后应注意三点：其一，未有十分把握之时，则守为主，取内线作战为研究要点。但一有机缘，则直取攻势，不可使其逝去也。此匪丑今日取攻势之行动，但其败兆亦即基于此，可以必也。其二，应力避正面一线配备，而转取纵深集团配备。只要吾人能坚持固守，则匪部交通接济，均无策源，必难持久，此亦必然之势也。其三，则游击战争与伏兵战争，急须实施讲求，并设法奖勉，否则无异盲目失聪，终为匪束手宰割也。\footnote{《陈诚先生回忆录·国共战争》，台北“国史馆”，2005，第23页。}
\end{quote}


1933年12月，在第五次“围剿”初期作战告一段落后，蒋总结战争经验，认为：“土匪来袭击我们，差不多每次都是下午，尤其是下午六点钟以后……我们对付他最稳当的办法，就是尽早到达宿营地切实准备，使他无隙可乘，不敢来犯。”因此他规定国民党军行军作战两原则：“以后在匪区行军，每日以三十里至四十里为原则”；“下午二时以前全部到达预定的宿营地点，迅即集结宿营，二时以后，不得继续行军”。\footnote{蒋介石：《为闽变对讨逆军训话——说明讨逆剿匪致胜的要诀》，《先总统蒋公思想言论总集》第11卷，第633～635页。}应该说，蒋介石一再耳提面命的上述战略战术理论上并无过人之处，他本人也明确谈道：“我们现在以剿匪所最用得着的，乃是十八世纪末与十九世纪初期拿破仑战争时代所用的战术。”\footnote{蒋介石：《军官训练团训练的要旨和训练的方法》，《庐山训练集》，第19页。}但正如他所判断的：“匪区纵横不过五百方里。如我军每日能进展二里，则不到一年，可以完全占领匪区”。\footnote{蒋介石1933年10月17日战字二一三号训令，《赣粤闽湘鄂北路剿匪军第三路军五次进剿战史》（上），第三章，第12页。}战斗毕竟是以战胜为目的，由于国民党政权掌握着物质和人力资源的绝对优势，加上内外环境给了其从容展开的时间，稳扎稳打，步步推进，虽然拙笨，但却最有成算。


为全力限制红军机动，尽力压缩红军的活动空间，使红军擅长的机动作战方式难以发挥，除进军时注意稳扎稳打外，国民党军采用碉堡战术，通过大量构筑碉堡，对红军活动区域实施封锁并截断红军的活动路线。据碉堡战术的最早建议及采用者之一戴岳回忆，早在1929年冬江西召开的全省“清剿”会议上，第十二师师长金汉鼎发言时即提到：

\begin{quote}

清政府在云南、贵州等地镇压土著居民的反抗时，土著民族曾依靠建碉守卡的办法予清军以重大打击；后来清军也是采用这个办法，最后征服了土著民族的顽强抵抗；因而主张对红军的“进剿”也可仿效这个办法，借以巩固“进剿”部队的阵地，并进而逐步压缩苏区，最后“消灭”红军和革命根据地。\footnote{戴岳：《我对蒋介石建议碉堡政策的经过》，《文史资料选辑》第45辑，文史资料出版社，1964，第172页。}
\end{quote}


这一建议当时虽颇引起共鸣，但未付诸实施。1930年戴岳奉命“清剿”赣东北苏区时，对碉堡战术作了尝试，取得一定效果。1931年初，他根据这一经验提出“实行碉楼政策”，即扼绝要道，实行清乡，“断绝匪区的接济和通讯”的“围剿”方略。具体解释为：“碉楼的利处，就是能以少数的兵保守一方，使匪共不能击破，并且能以少数的部队击溃多数的股匪，同时可以阻绝匪的交通及活动，逐渐把匪区缩小。”\footnote{戴岳：《对于剿匪清乡的一点贡献》，转见冰澈《对于白军“剿共”的研究》，附录，第6、10页。}戴岳的这一主张，当时由于环境及认识的限制，仍未得到普遍推广，但重要据点碉堡的修筑已陆续开始。1932年6月，蒋介石在日记中写道：“阅三省剿匪计划，修正之。此次剿匪经费决以半数为修路、筑碉、赈济之用，而治本之道则注重于清廉县长与组织保甲、训练民团，分配土地，施放种籽也。决不求其速剿，只望其渐清也。”\footnote{《蒋介石日记》，1932年6月11日。}修路、筑碉的作战思路在蒋这里引起注目。第五次“围剿”准备工作开始时，蒋介石即坚决贯彻大量构筑碉堡、利用碉堡层层推进的碉堡政策，强调：“清剿各部到处以修碉筑堡为惟一要务。”\footnote{《蒋中正电清剿各部到处以修筑碉堡为要务（1933年7月29日）》，蒋中正文物档案002020200019065。}对碉堡的运用，他在下面一段话中说得很清楚：

\begin{quote}

从前曾国藩李鸿章他们在淮水及黄河流域清剿捻匪，曾掘几千里的长沟，其工程之浩大，决非我们现在掘几里几十里战壕可比。但是壕沟在北方原很相宜，到了南方多山的地方，便不适用，所以我在江西剿匪，改用碉堡，在两个碉堡之间，我们的火力可以交叉射击，敌人便不能通过。这样连下去由福建经江西连到湘鄂，差不多有一二千里之长，只要每里以内平均有两个碉堡，敌人便不容易窜入。所以碉堡线可说是我们的万里长城，这个长城筑得坚固，就可以困死土匪！\footnote{蒋介石：《剿匪与整军之要道》，《先总统蒋公思想言论总集》第13卷，第221页。}
\end{quote}


碉堡战术确定后，蒋介石对碉堡建设十分重视，规定“每星期一连必须添筑一个以上之碉堡”。\footnote{《蒋介石巧戌行战六手令（1933年11月18日）》，《赣粤闽湘鄂北路剿匪军第三路军五次进剿战史》（上），第四章，第21页。}部分碉堡可由民众帮助守卫，“分给其手榴弹与旧枪子弹，使之不赖军队而能于一个月内切实自卫也”。\footnote{《蒋中正电赵观涛实地查察信河北岸修筑碉堡等事》，蒋中正文物档案002010200090044。}第五次“围剿”初期，国民党军有些部队对构筑碉堡执行不力，1933年11月中旬，红一军团为策应三军团发动的浒湾攻击战，北上突破国民党军在见贤桥、永兴桥一带构筑的松散碉堡线，令蒋介石为之震动。他连电前方，要求改进碉堡配置，增加密度，“封锁碉堡群之间隔，不得过二里以上”，同时指示加强碉堡群之构筑，形成重叠立体配置和相互间的有效火力配合。蒋介石并严厉警告，如构筑碉堡“再因循玩忽，查出定以军法从事，决不姑宽”。\footnote{《蒋介石皓酉行战六手令（1933年11月19日）》，《赣粤闽湘鄂北路剿匪军第三路军五次进剿战史》（上），第四章，第21页。}在蒋介石一再督促下，“敌土工作业力确实增强”。\footnote{《彭滕关于七日战况和改变行动的报告（1933年11月7日）》，《中央苏区第五次反“围剿”》（上），第100页。}国民党军第五次“围剿”期间的作战和推进始终与堡垒修筑同步，到1934年4月，仅在江西构筑各种碉堡5300余座，\footnote{《江西省保安处工作报告摘要》，《南昌行营召集第二次保安会议记录》，南昌行营编印，1934，第31页。}福建573座。\footnote{《福建省保安处工作报告摘要》，《南昌行营召集第二次保安会议记录》，第35页。}“围剿”终时，其主攻部队北路军第三路军构筑碉堡4244座，\footnote{《赣粤闽湘鄂北路剿匪军第三路军五次进剿战史》（下），附表：《本路军戡匪全战役中各部队构筑碉堡公路统计比较图》。}第六路军922座，\footnote{《赣粤闽湘鄂剿匪军北路第六路军赣南围剿之役构筑碉堡一览表》，《剿匪纪实》（上），第六路军总指挥部，1937，第35页。}为沟通碉群大力修筑公路，1932～1934年间，共修筑公路4297.22公里，相当于1931年江西公路总里程284.6公里的十多倍。\footnote{孙兆乾：《江西农业金融与地权异动之关系》，萧铮主编《民国二十年代中国大陆土地问题资料》第86种，第45236页。}碉堡之多，即便是江西省会南昌近郊，也是“堡垒林立”。\footnote{郑贞文：《赴赣日记》，《福建教育厅教育周刊》第192期，1934年6月4日。}


第五次“围剿”期间国民党军构筑的大量碉堡及其通连公路组成的严密封锁网，既可弥补国民党军战斗精神的不足，又针对着红军缺少重武器、难以攻坚的弱点，强己抑彼，一箭双雕，对红军机动作战造成相当困难。反“围剿”初期的硝石战役中，红军集中一、三、五、七军团7个师部队进攻硝石，期望打破敌军“围剿”兵力部署，但指挥作战的彭德怀回忆，由于“当时黎川驻敌三、四个师，南城、南丰各约三个师，硝石在这三点之间，各隔三、四十里，处在敌军堡垒群之中心。我转入敌堡垒群纵深之中，完全失去机动余地，几乎被敌歼灭”。\footnote{《彭德怀自述》，人民出版社，1981，第185页。}国民党方面战后总结该战役获胜主要原因也即“工事坚强”。\footnote{《赣粤闽湘鄂北路剿匪军第三路军五次进剿战史》（上），第二章，第27页。}林彪对碉堡战术也有很深印象，谈道：“敌人每到一地，他立即进行筑垒，以立定脚跟，接着构筑联络堡、封锁线、马路，以取得别的联络与策应……在前进中如遇到我军稍有力或有力的阻滞时，他立即停止向预定的目标前进，而进行筑堡，与防御的战斗。”“这些方法，都是着眼于使红军在政治军事上的优越条件困难充分运用，而使他自己在物质上兵器的优点能极力发扬”。\footnote{林彪：《短促突击论》，《革命与战争》第6期，1934年7月。}事实上，当时苏区内存在的“土围子”即是碉堡战术效用的一个有力证明。李德回忆，他在中央苏区指挥作战时，有一次经过“一个用岩石筑成的堡垒”，“堡垒的洞穴里驻有一个地主的民团，因为他们有充足的粮食储备和水源，所以能坚守至今”。\footnote{〔德〕奥托·布劳恩：《中国纪事1932～1939》，现代史料编刊社，1980，第61页。}这样的“土围子”在苏区并不鲜见，有的坚持达三四年之久。单独的“土围子”在苏区范围内的坚持，和红军缺少重武器而难以攻坚直接相关。


和碉堡战术相应，在“围剿”进行过程中，蒋介石具体作战指挥时坚决贯彻稳中求进方针。第五次反“围剿”初期，红军曾集中主力部队出击赣东资溪桥、硝石一带，以红五军团牵动敌人，红一、三军团准备在敌被牵动时实施突击，这是红军擅长的一贯打法。但国民党军根据蒋介石的战略安排，“十分谨慎，步步为营，稳扎稳进，很少出击”，结果红军既“未能牵动敌人”，\footnote{《李志民回忆录》，解放军出版社，1993，第227页。}投入攻坚战后又无成算，不得不撤出战场。广昌战役初期，蒋介石明确指示陈诚：“我中路军主力不必求急进，只要固守现地，作成持久之局，以求薛、汤两路之发展，则匪经此战必崩溃更速，不必心急也。”\footnote{《手谕枫林三溪一线为我军进入赣南之第一步并示作战要领（1934年3月13日）》，《陈诚先生书信集·与蒋中正先生往来函电》（上），台北，“国史馆”，2007，第131页。}在进攻已有进展时，他又一再指示部队应继续辅之以碉堡和公路线之推进，“进展不必过急”，\footnote{《蒋介石1934年5月1日致蒋鼎文电》，《中华民国重要史料初编——对日抗战时期》绪编（二），第403页。}“总须稳扎稳打为要”。\footnote{《蒋介石1934年5月1日致汤恩伯电》，中国第二历史档案馆编《中华民国史档案资料汇编》第5辑第1编军事（4），江苏古籍出版社，1994，第322页。}陈诚在攻下广昌后也谈道：

\begin{quote}

近来听得许多官兵说：“土匪已经不得了了，我们可以直捣土匪老巢而剿灭之”。这种意思，固然很好，但是我们要知道，我们的剿匪要策出万全才行。现在步步筑碉的办法，是最稳是最妥的方法……我们一步一步的前进，土匪灭亡，是可计日而待的；不要性急，小心谨慎，才是要着。\footnote{《总指挥训话记录》，《汗血月刊》第1卷第5期，1934年7月20日。}
\end{quote}


1934年9月，整个“围剿”战局已有尘埃落定之势时，蒋接获前方进军计划后，仍强调部队“不宜单独……进展，应令会同各纵队前进，免为匪乘”，\footnote{《蒋介石1934年9月19日致顾祝同电》，《中华民国史档案资料汇编》第5辑第1编军事（4），第109页。}“仍宜注意集结，勿过分分散”。\footnote{《蒋介石1934年10月11日电》，《中华民国史档案资料汇编》第5辑第1编军事（4），第269页。}这体现出蒋对“围剿”作战极端慎重、务求必胜的态度。在总体力量占据绝对优势的情况下，承认自身与红军战斗精神和战斗能力的差距，将自己客观地摆到弱势地位，以拼红军的姿态确定战略战术方针并指挥作战，这对心高气傲的蒋介石虽然不一定那么情愿，却是他摆脱前几次“围剿”被动局面，在双方战略对抗中抢得先机的关键一步。

\section{国民党军作战基础的增进}


蒋介石在第五次“围剿”中选择持久消耗作战方针，和当时国内外相对有利的环境直接相关。周恩来曾经谈道，蒋介石在“第五次‘围剿’中，能出动五十万大军发起进攻和实行封锁，那是他的全盛时期”。\footnote{〔美〕艾德加·斯诺：《红色中华散记（1936～1945）》，群众出版社，1991，第68页。}确实，和前四次“围剿”几乎一直在国内外动荡局势中进行相比，第五次“围剿”进行过程中，南京政府内外环境相对宽松，给了其从容展布的空间。


从外部环境看，长城抗战并签订《塘沽协定》后，日本在华北的侵略活动告一段落，北方的压力暂时有所减轻，此后直到第五次“围剿”结束，日本在华北一直未有大的动作，南京政府获得相对稳定的外部环境。与此同时，南京政府积极调整对外政策，与英、美等国加强联系，行政院副院长、财政部长宋子文于1933年4月开始长达半年的欧美之行，并与美国订立5000万美元的棉麦借款合同。宋子文之行被认为标志着“南京政府对欧美国家实行经济开放政策的起端”。\footnote{石源华：《中华民国外交史》，上海人民出版社，1994，第410页。}在加强经济联系同时，南京政府向西方国家大量订购武器装备，据中央信托局统计，1933年和1934年两年间，购买军火费用达6000多万元。\footnote{《中央信托局经办各项军械军火及航空器材数额统计图（1936年）》，中国第二历史档案馆藏。}这些，既加强了南京政府与西方国家间的政治、经济联系，又提高了其军事装备和统治能力。共产国际军事顾问1933年初报告：“在敌人方面，武器装备的质量在不断提高，其机枪和火炮的配备程度已接近现代军队的水平。”\footnote{《布劳恩关于中央苏区军事形势的书面报告（1933年3月5日）》，《共产国际、联共（布）与中国革命档案资料丛书》第13册，第336页。}具体到实战中，国民党军发动第五次“围剿”后，将其装甲部队和新购德国山炮投入进攻，发挥出相当效果。粟裕回忆：“十九师是红七军团的主力，战斗力强，擅长打野战，但没有见到过装甲车……部队一见到两个铁家伙打着机枪冲过来，就手足无措，一个师的阵地硬是被两辆装甲车冲垮。”\footnote{《粟裕战争回忆录》，解放军出版社，1988，第103～104页。}红军两个主要军团的指挥者彭德怀和林彪都注意到：“蒋军在第五次‘围剿’时，技术装备比以往几次有所加强。”\footnote{《彭德怀自述》，第188页。}“每连有多至六挺的机关枪，至少也有一挺。我们在敌机枪下除非不接近，一接近一冲就是伤亡一大堆。”\footnote{林彪：《关于作战指挥和战略战术问题给军委的信（1934年4月3日）》，转见黄少群《中区风云——中央苏区第一至第五次反“围剿”战争史》，中共党史出版社，1991，第324页。}周恩来在1934年初对反“围剿”战争作出初步总结时，印象很深的就是国民党军装备和编制的加强：

\begin{quote}

每团中的火力配备改变了。过去每营有机关枪连，现改为每团属一重机关枪连，而每一步兵连中则加多轻机关枪的数目（每连三架）。

他的部队中的非战斗人员亦大大裁减，为便于山地运动与作战。他的行军公文也减少到最小限度，增加战斗人员的比例。\footnote{周恩来：《五次战役中我们的胜利（一）——论持久战》，《红星》第33期，1934年3月18日。}
\end{quote}


随着对国内地方实力派的压迫、清除和国内建设的展开，南京政府统治力量也在逐渐增强，内部形势趋于稳定。第二、三、四次“围剿”中，南京政府均遇到严重的内部纷争或外敌压迫。中共有关报告显示，第二次“围剿”时，国民党军起初也取稳扎稳打方针，后因胡汉民事件导致粤变不得不加快“围剿”行动，使红军觅得作战良机。尤其当红军首战东固时，蔡廷锴部在红军后方，红军仅有一师部队牵制，“假使此时蒋蔡北进，则我们后方受到危险非常困难，但此时蒋蔡不但未前进，并且将前进到城冈的部队撤退到兴国城（原因大概是因为广东问题发生）”，\footnote{《中央苏维埃区域报告》，《中央革命根据地史料选编》（上），江西人民出版社，1982，第370页。蔡廷锴本人也回忆，1931年5月粤方公开反蒋后，他们作为与粤、蒋双方均有关系部队，处境两难，“心灰意冷，对于剿赤任务，亦只得放弃。我与蒋总指挥即决心回师赣州，静观时局之演变”。见《蔡廷锴自传》（上），黑龙江人民出版社，1982，第241页。}使红军摆脱了腹背受敌的危险。第三、四次“围剿”期间的粤桂反蒋、九一八事变、长城事变，也直接影响到国民党军的作战行动。第五次“围剿”中，虽然1933年底发生以十九路军为军事基干的福建事变，但由于事变很快遭镇压，蒋介石反而由此取得对福建的完全控制。福建事变的迅速失败，也使各地反蒋力量不得不有所忌惮，此后相当一段时间内，地方实力派的反蒋公开活动偃旗息鼓。完全控制福建后，除广东陈济棠外，国民政府从东、西、北三面完成对中央苏区的包围，无论是军事展开或经济封锁都获得了更好的内外环境。尤其是鄂豫皖、湘鄂西苏区红军主力相继转移后，长江中游地区基本安定，蒋介石更可全力筹划对付中央苏区之策。1933年6月23日，蒋介石曾在日记中写道：“剿匪根本计划渐立，逐渐推进，期以一年，必能成功，惟时局是否许可耳。”\footnote{《蒋介石日记》，1933年6月23日。}后来的事实证明，蒋介石获得了这种时机。直接在前线指挥“剿共”军事的陈诚对此颇有同感，第五次“围剿”发动之初，他在家书中表示：“此次剿匪当有把握，因蒋先生亲自督剿，无论精神物质方面，均有关系。以部队数量言，已较之以前增加十余师。而精神方面，各级经庐山训练之后，确有进步也。”\footnote{《蒋先生亲自督剿此次剿匪当有把握（1933年10月13日）》，《陈诚先生书信集·家书》（上），第235页。}


经过庐山训练及部队整训，国民党军队整体战斗力和战斗精神确有一定提高。第五次“围剿”期间，国民党军队中高级官员大都能身先士卒，出现在战场第一线。从主攻部队第三路军伤亡情况看，该路军伤亡军官总数988人，伤亡士兵总数10755人，\footnote{《赣粤闽湘鄂北路剿匪军第三路军五次进剿战史》（下），附表，第三路军五次进剿全战役官兵伤亡统计总表。}军官和士兵之比约1∶11，高于战斗部队实际官兵比。军官伤亡较重和其在第一线指挥作战相关。第五次“围剿”时，国民党军师、旅长一级指挥员经常出现在战场第一线，团长一级指挥员则直接指挥并于必要时参加战斗。德胜关、寨头隘、大罗山等战役中，第六十七师三九九团团长、七十九师二三五旅四七四团团长、六师十八旅三十六团团长都能身先士卒，亲率所部与红军展开肉搏战，第六师十八旅三十六团团长并当场阵亡。国民党方面总结战役得手主要原因即为：“各高级指挥官身临前线，从容指挥。”\footnote{《赣粤闽湘鄂北路剿匪军第三路军五次进剿战史》（上），第五章，第16页。}凤翔峰一役，四七〇团第二、三营代营长均先后伤亡，“连长以下干部伤亡已达二十余员”，“第六连之连排长伤亡殆尽，仅赖一军士毛炳芳指挥”，\footnote{《赣粤闽湘鄂北路剿匪军第三路军五次进剿战史》（上），第六章，第23页。}但仍能守住阵地。鸡公山战斗中国民党军也异常奋勇，红军战报记载，当时守军为红军一个营部队，在团长、营长率领下坚守，“当敌迫至工事近前时，该营长即令战斗员将炸弹收集起来连掷五六弹，伤敌三十余名，但敌仍继续攻击，该营长旋即牺牲，不久该团长亦受伤，而部队失了掌握，加以手榴弹告绝，子弹所有亦无几系（原文如此——引者注），致放弃阵地”。\footnote{《九军团鸡公山战斗经过详报（1934年2月9日）》，《中央苏区第五次反“围剿”》（下），第57页。}国民党军的上述表现，为红军指挥员所注意，时任红五军团十三师政治部主任的莫文骅回忆，第五次反“围剿”期间，国民党军表现颇为顽强，“尽管敌人在红军阵地前倒下了一大片，但后面敌人还是一股劲往前冲”。\footnote{《莫文骅回忆录》，解放军出版社，1996，第230页。}整个“围剿”期间，国民党军“几乎没有起义者，只有很少的被俘者”。\footnote{〔德〕奥托·布劳恩：《中国纪事1932～1939》，第55页。}


注意发挥前线指挥官主动性，不过多干预具体作战，是蒋介石在第五次“围剿”期间指挥作战的一个突出特点，而前线指挥官也表现出相当强的主动精神，对战役胜利发挥了重要作用。1934年1、2月间，陈诚与蒋介石曾围绕着福建事变后的主攻方向展开争论，最后以蒋介石的退让而告终。1934年2月，汤恩伯第十纵队和刘和鼎第九纵队准备进攻闽西北的将乐、沙县，然后与第三路军配合进窥建宁。17日，蒋介石致电刘、汤，令其迅速在将乐、沙县发动，要求其“就近审度情势，从速断行，但无论如何，将乐应派队先行占领，俾得相机截击匪部西窜之路”。\footnote{《蒋介石1934年2月17日致刘和鼎、汤恩伯电》，《中华民国史档案资料汇编》第5辑第1编军事（4），第293页。}18日，刘、汤致电蒋介石，同意“以主力使用于将乐方面”。\footnote{《刘和鼎、汤恩伯1934年2月18日致蒋介石电》，《中华民国史档案资料汇编》第5辑第1编军事（4），第293页。}22日，鉴于将乐处沙县、建宁之间，先攻将乐，有“两侧均受威胁”\footnote{《东路军第十纵队与红军在闽北一带战斗详报》，《中华民国史档案资料汇编》第5辑第1编军事（4），第294页。}之虞，刘、汤决定改变计划，先攻沙县，并立即付诸实施。2月25日和3月6日，刘、汤所部先后进占沙县、将乐。根据以占领要地而不以实施歼灭战为主的“围剿”总方针，刘、汤先攻沙县的行动有效却不失稳重，而蒋实施中间突破、先占将乐切断建宁红军退路的想法虽然凶狠，但风险相对较大，和稳中求进的总思路不无抵触。正因如此，虽然刘、汤的行动和蒋的具体作战指示相违背，蒋事后却未追究。对此，东路军有关战报总结为：“匪踪飘忽无常，情况变化靡定，关于各部队作战部署与进展步骤为适时到达，皆以电达要旨命令，俾前方指挥官，得以酌量情况，敏活运用。”\footnote{《国民党东路军蒋鼎文部与红军在闽赣境内战斗详报》，《中华民国史档案资料汇编》第5辑第1编军事（4），第165页。}得意之色溢于言表。第五次“围剿”中，陈诚、汤恩伯等国民党军前线将领在战役指挥上的机动处置及蒋介石对其意见的尊重和接受，既调动了前线将领的积极性和主动精神，也使国民党军作战指挥和作战方针更切合前线实际。

\section{国民党军的作战部署}


为全力准备发动第五次“围剿”，蒋介石在军队的组织、编制和调配上作了一系列部署，殚精竭虑，全力以赴。1933年5月，蒋介石决定成立军事委员会委员长南昌行营，\footnote{1931年2月，蒋介石曾设立“陆海空军总司令南昌行营”。1933年初第四次“围剿”时设立军事委员会委员长南昌行营。}委任江西省主席熊式辉兼行营办公厅主任；行营原参谋长贺国光为第一厅厅长，主管军事；行营秘书长杨永泰兼第二厅厅长，主管政治。行营全权管理赣、粤、闽、湘、鄂五省党军政要务。并由南昌行营组设党政军设计委员会，为五省一切党政军事务参谋部。随后，蒋介石坐镇南昌，亲自指挥部署对苏区的第五次“围剿”。为使部队编制更为精干，适应对红军山岳地带作战的特点，南昌行营决定改变部分军队的编制，将每师辖两旅、每旅辖两团或三团的部队改编为每师直辖三团，增加战斗人员编制。团以上机关增加侦察队和扩大运输队，以灵敏耳目、保证给养。


具体改编结果是：第四军九〇师改编为五十九、九〇两个师；第十八军一一师改编为一一、六十七两个师，十四师改编为十四、九十四两个师，四十三师改编为四十三、九十七两个师；第三十六军第五师改编为五、九十六两个师，独立三十二旅、三十三旅改编为九十二师、九十三师。第九十八师及第五十二、五十九两师合编的第九十九师仍保留两旅六团建制。


第五次“围剿”开始时，国民党军作战的具体部署是：以驻赣、粤、闽、湘、鄂各省部队分编为北路、南路和西路军。北路军以顾祝同为总司令，蒋鼎文为前线总指挥，下辖3路军，共33个师另3个旅，是“围剿”中央苏区的主力。其编成和任务是：以4个师、1个旅及第二纵队、税警总团编为第一路军，顾祝同兼总指挥，刘兴为副总指挥，配置于新干、吉水、永丰、乐安、宜黄地区，在此构筑碉堡封锁线，逐步向中央苏区推进，并阻止红军向赣西北前进；以6个师编为第二路军，蒋鼎文兼总指挥，汤恩伯为副总指挥，配置于祟仁、藤桥、金溪地区，构筑碉堡封锁线，逐步向中央苏区推进，并阻止红军向赣东北前进；以18个师另1个补充旅编为第三路军，陈诚为总指挥，薛岳为副总指挥。该路军又以14个师编为机动作战的第五、第七、第八纵队，以4个师1个旅编为守备队，集结于南城、南丰地区，沿抚河两岸构筑碉堡封锁线。第三路军是北路军中的主力，其任务是：在第一、第二路军的配合下，依托碉堡向广昌方向推进，寻求红一方面军主力决战。此外，北路军总司令部还直接指挥第二十三、第二十八师，扼守赣江西岸的吉安、泰和等地，配合西路军封锁赣江，阻止红一方面军主力西进。第十三、第三十六、第八十五师为总预备队，由钱大钧统领，置于抚州附近地区。浙赣闽边区警备部队5个师另4个保安团由赵观涛统领，“围剿”闽浙赣苏区，并配合北路军第二路军阻止红军向赣东北方向发展。


西路军以何键为总司令，指挥14个师又2个旅，分三个纵队。第一纵队配置于湘赣边境酃县、茶陵、莲花地区，负责“进剿”永丰、宁冈一带红军；第二纵队配置于浏阳、平江、铜古地区，负责“进剿”湘鄂赣边红军；第三纵队配置于大冶、通山地区，协同第二纵队“围剿”鄂南边境之红军。


南路军以陈济棠为总司令，指挥粤军11个师又1个旅，筑碉扼守武平、安远、赣县、信丰、上犹、崇义地区，阻止红军向南发展，井逐步向筠门岭、会昌地区推进，协同北路军作战。


福建方面第十九路军等部共6个师又2个旅，扼守闽西和闽西北地区，阻止红军向东发展。


值得注意的是，国民党军的部署中包含有大量守备部队，这使国民党军配备更为纵深，后方防御力量大为加强，这也是其碉堡战术的必要配套措施。对此，共产国际军事顾问谈道：“敌人将其兵力划分为守备部队和突击部队。这样，虽然一方面他们的作战兵力削减了50\%多，但另一方面，我们几乎不可能将自己对敌军的战术胜利扩大到真正意义上的战略胜利，因为红军不得不一次又一次地面对有足够数量的新敌军保卫的防御工事。”\footnote{《布劳恩关于中央苏区军事形势的书面报告（1933年3月5日）》，《共产国际、联共（布）与中国革命档案资料丛书》第13册，第340页。}


第五次“围剿”初期，国民党军用于“围剿”的空军共有五队，拥有飞机50架，绝大部分是侦察机（包括可塞、菲亚特、道格拉斯等机型）兼用于轰炸，马力、飞程、载重力都十分有限，部署于南昌、抚州、南城等地。后期，南京政府又订购新式战机投入战场。李德当时报告：“敌人现正使用新的空战有威力的兵器，并在空中用新的技术……引用新式有力的轰炸机（六百匹马力不论天时气候都能飞，并且有很大的载重力和飞程）。”\footnote{华夫（李德）：《论防空几个新的问题》，《革命与战争》第8期，1934年8月15日。}李德所提到的这批新飞机，国民党军战史有清楚的记载：“是年9月起，新机陆续抵达。先分配空军第三、四、五三队，每队‘新可塞’机二架。十一月，拨交第一队美制‘诺斯罗卜’轰炸机九架，第二队‘新可塞’机九架。”\footnote{《国民革命军战史》第二部《安内与攘外》（一），台北，“国防部史政编译局”，1993，第634页。}这时，第五次“围剿”已处于尾声阶段。


国民党方面的第五次“围剿”，既占据着中央政府的制高点和全国资源的绝对优势，又利用着国际国内的有利形势，大军毕集，志在必得，给中央苏区造成了空前压力。

\section{“七分政治”的具体实施}

\subsection{蒋介石的“三分军事，七分政治”}


第五次“围剿”期间，吸取前几次“围剿”失败的教训，蒋介石依据“三分军事、七分政治”原则，将“围剿”定位为军事、政治、经济、社会的总体战，这一军事、政治并用方针，对“围剿”的最终走向发挥了一定作用。


作为一个国民革命时期曾经在政治宣传、鼓动上大得人心并获取力量的政党，国民党深知政治宣传、收揽人心的重要。在与中共的对垒中，国民党方面政治上虽然不像中共那样游刃有余，握有主动，但一直力图有所作为。早在1931年，何应钦就提出：“要消灭共匪，非党政军全体总动员集中力量团结意志不能挽救危机，军事只可以治标，正本清源以及休养生息的种种任务，是望政府和党部来担当责任。”\footnote{《在赣欢宴各界之演说》，1931年2月26日《江西民国日报》。}蒋本人也谈道：“本总司令于去岁督师江西之时，即深知剿灭共匪与寻常对敌作战绝对不同，苟非于军事之外同时整理地方，革新行政，断难以安抚民物而奏肃清之功。”\footnote{《蒋中正呈送剿匪区内各省行政督察专员公署组织条例（1932年8月29日）》，台北“国史馆”藏国民政府档案，001012071114。}这样的反省尚称深切，已切实意识到双方胜败的关键所在。


从1931年开始，南京政府采取了一系列措施，力图从社会政治等层面强化本身力量，以与中共强有力的政治组织抗衡。南昌陆海空军总司令行营成立江西地方整理委员会，整顿江西地方政治、社会。主要措施包括：督促整顿全省保卫团，将各县反共义勇队一律改组为保卫团，区团以下均设守望队，形成全面监视巡查网；蠲免1930年度全省田赋、地租；制定《处理被匪侵占财产办法》，规定赤化收复区域土地、房屋各归原主，恢复地方秩序；颁布《保护佃农暂行办法》，规定地租最高额不得超过百分之四十，如遇天灾，佃农要求减租，地主不得拒绝；组织由地方逃至中心城市的“难民”中的青壮年随军返乡，协助运输、带路，或参加筑路。1931年初，在南京中央政治学校开办特别训练班，下分“剿匪宣传队”，施以政治宣传训练，再以团为单位分配到前方部队，指导政治和宣传工作。在对特别训练班的训话中，蒋介石数次提到“剿匪的实施宣传要占六分力量，军事只能占四分力量”，\footnote{《蒋中正总统档案·事略稿本》第10册，第50页。}这应可视为其后来提出的“三分军事，七分政治”的原始版。6月，南京中央发布对各级党部训令，要求加强与中共全方位的政治争夺，特别指示“组织健全的巡回乡村宣传队”，进行“剿匪宣传”。\footnote{《蒋中正总统档案·事略稿本》第11册，台北，“国史馆”，2004，第282～283页。}为配合“剿共”军事、增进行政效率，南昌行营设置“党政委员会”，蒋自兼委员长，将江西全省“剿匪”区域共43县划分为9个分区，每区设置党政委员会分会，负责指导各区军事、政治、经济等事务。\footnote{《党政委员会组织成立呈报中央文（1931年7月13日）》，蒋中正文物档案002060500009012。}这些措施虽由于缺乏具体组织实施的决心和能力，实际效果有限，如蒋介石自己所说，“劳师转馕，苦战连年，地方贤良士民，竟无出而相助者”，\footnote{《蒋委员长告江西各县离乡避匪之贤良士民书（1933年9月22日）》，《军政旬刊》第2期，1933年10月30日。}但起码表明国民党方面对自身在地方政治建设上的薄弱环节已有所注意，开始努力在政治和组织上与中共争夺民众。


在参加“剿共”战争的过程中，一些国民党将领也意识到与中共在政治上展开争夺的重要性。作为“剿共”前线指挥官，曾任国民党军第十八师第五十二旅旅长的戴岳深刻体会到民心向背对武力利钝的影响，明确指出：“清剿匪共，绝不是军队一部分的力量做得到的，是要党、政、军、民通力合作才行的。”\footnote{戴岳：《对于剿匪清乡的一点贡献》，转见冰澈《对于白军“剿共”的研究》，“附录”，第11页。}运用政治力量，在政治上与中共展开争夺，首先就必须取得民众的支持，将民众从中共方面拉到自己一边，这是政治战的基本。为此，他不赞成何应钦等提出的“进剿”部队“除班长及由官长指定之士兵外，概禁止与人民接谈”、\footnote{何应钦：《剿共须知》，转见冰澈《对于白军“剿共”的研究》，附录，第32页。}“禁止士兵无故与民众往还”\footnote{郭汝栋：《国民革命军第廿军剿匪部队注意事项》，转见冰澈《对于白军“剿共”的研究》，“附录”，第81页。}的主张，而强调“所到之处，要随时召集民众开会，揭破共匪的阴谋，宣扬本党的三民主义”。\footnote{戴岳：《对于剿匪清乡的一点贡献》，转见冰澈《对于白军“剿共”的研究》，“附录”，第19页。}从争取民众的目标出发，他们强烈反对杀戮苏区民众，南昌行营第二科科长柳维垣列举了烧杀政策的危害：“匪屋不烧，或尚有悔过反正之日，一烧其屋，即迫其终身从匪……正合共党之希望。”\footnote{柳维垣：《剿匪实验见解录》，转见冰澈《对于白军“剿共”的研究》，“附录”，第41页。}


为改变国民党政权在民众中的糟糕形象，他们还颇有眼光地提出对地主、土豪的态度问题。作为中共革命的打击对象，地主、土豪往往投向国民党方面，成为国民党军的主要依靠对象。邻接中央苏区的湘赣省曾对地主、豪绅和反革命及其家属实行驱逐政策，结果这些人“驱逐到白区以后，就参加国民党军队，带领地主武装和国民党的部队来打我们。他们地形、路线很熟，而且大多是青年，一出去反动得很，对我们很凶，使我们吃亏不小”。\footnote{张启龙：《湘赣苏区的革命斗争》，《湘赣革命根据地》（下），中共党史资料出版社，1991，第863页。}不过，国民党军依靠这些人固然可以收到一定效果，但在苏维埃革命的风潮中，完全站在地主豪绅一边就意味着和大多数的普通农民对立，这对国民党军争取更多民众的支持十分不利。所以戴岳特别提醒：“难民是各村逃出来的人，对于地形道路匪情都是很熟悉的，可以把他组成梭镖队，随军带路；但是他们没有纪律，加以报仇心切，所到之处，随意烧杀，这是要特别注意纠正的。”要求遏止“还乡团”的疯狂报复行为，“设法和解难民向反共来归的农民寻求报复”，\footnote{戴岳：《对于剿匪清乡的一点贡献》，转见冰澈《对于白军“剿共”的研究》，“附录”，第11～12、20页。}希望通过与土劣保持距离，改变国民党在普通农民中的富人利益维护者形象，尽可能争取更多人民支持。


戴岳等国民党军将领在与中共交战过程中，汲取经验教训，对前线实况，对苏区政治、社会、军事状况有更多的了解，他们的主张也逐渐为蒋介石所注意，戴岳、柳维垣都曾接受蒋的召见、垂询，其中不少意见为蒋所采纳。正是在此背景下，第五次“围剿”准备过程中，蒋介石把政治力量提到空前的高度，强调：“剿匪乃争民之战，非争地之战，故军事纵告胜利，如无健全之政治设施，相辅而行，则终必徒劳无功。”\footnote{《为剿匪军事行动之后须继之以政治设施财政支援事致汪兆铭院长电（1933年8月21日）》，《先总统蒋公思想言论总集》第37卷，第86页。}将政治争夺战置于军事之上，要求所部主动出击，与中共展开政治争夺：“政治工作人员之工作必须向匪区设法深入为唯一任务。”\footnote{《蒋中正指示政治工作须深入匪区（1933年10月27日）》，蒋中正文物档案002010200096040。}


和中共相比，国民党的鼓动性和组织力自是望尘莫及，这和两党的理论基础、奋斗目标、人员构成、领袖特质等诸多因素相关，非短期所能改变。基于对自身特点的了解，蒋介石强调政治的争夺不应好高骛远，而要注意于一时一地一事的实际解决，正如他此时谈到的：

\begin{quote}

一谈到经济设施，开口便说要如何统制，这些都是不切实际的理想，亦就是没有用的理论，都不是目前我们所需要的。我们现在最需要的，就是要注意研究一切眼前的实际问题，完全针对客观的事实，一件一件从实地调查考察来拟订具体能行的解决办法。我们现在更不可憧憬于什么高远的理想，亦不必发表新奇的理论，我们只是竭忠尽智为国计民生来打算，就事实来求解决，从现实的工作中来求进步……如果我们能够将这些事情，一件一件的改进，将大大小小的事实问题，一个一个的解决，已经够了。\footnote{蒋介石：《革命成败的机势和建设工作的方法》，《先总统蒋公思想言论总集》第11卷，第606～607页。}
\end{quote}


本着这一认识，第五次“围剿”期间，国民党在政治、经济、社会组织等方面全盘进行整合，推出一系列具体措施。


其一，对苏区民众和红军展开攻心战术，改变国民党政权和国民党军的负面形象。南京国民政府成立后，由于基层力量薄弱，各地乡绅多被作为当地社会代表，负上传下达之责，成为政府控制基层社会的重要依靠力量。其中不良分子往往利用权势徇私舞弊，坑害百姓。为改变国民党政权的富人维护者形象，“围剿”数遭失败后，蒋介石和国民党政权开始反思此前对乡绅的依赖政策，希望与乡村中的权贵阶层保持距离，限制土劣活动。第五次“围剿”前，国民党方面着力调整其乡村政策。1933年4月，蒋介石通电各省政府，指出：“绅士仍多狐假虎威，欺下罔上之事。各区对于绅士固应多方物色吸引，但主管官应严加监察，推行政治，勿使阻隔。”\footnote{《电杨永泰令各省政府主管严加监督乡绅言行（1933年4月6日）》，蒋中正文物档案002010200081020。}对前方官兵则要求：“一定要亲近醇厚可用的真正的民众，尤其是一般真正的民众的领袖，决不好亲近一般土豪劣绅。”\footnote{蒋介石：《推进剿匪区域政治工作的要点》，《先总统蒋公思想言论总集》第11卷，第239页。}1933年8月，南昌行营在前一年鄂豫皖三省“剿匪”总司令部有关条例基础上，制定颁布《惩治土豪劣绅条例》，规定：“武断乡曲，虐待平民，致死或笃疾者处死刑，或无期徒刑”；“恃豪怙势，蒙蔽官厅，或变乱是非，胁迫官吏……者，处五年以上，七年以下有期徒刑”。\footnote{《国民政府军事委员会委员长南昌行营惩治土豪劣绅条例》，《国民政府军事委员会委员长南昌行营处理剿匪省份政治工作报告》，第十章，南昌行营编印，1934，第2页。}专门针对豪绅在法律框架内出台相关打击条例，体现出蒋介石和国民党政权与豪绅划清界限的愿望，对豪绅、地方官员及普通民众都是一个表态。根据这一条例，到1934年年中，何键主持的西路军共接办土劣案件48件，结案35件，其中相当部分都是针对武断乡曲、欺压百姓作出的判决。\footnote{《党政处报告关于清乡善后整训团队编查保甲户口等事项工作报告表》，《西路军公报》第12期，1934年6月15日，报告类第6页。}湖北江陵通缉巨绅周瑞卿，一度使附近“大小土劣相率敛迹”。\footnote{湖北省民政厅：《湖北县政概况》，台北，文海出版社，1976，第912页。}


整顿军纪是国民党军实行自我改造的重要一环。国民党军由于历史、现实的原因，有所谓中央军、杂牌军之分，中央军供应充足，军纪一般相对较好，杂牌部队则在供应和军纪上都难以保证。为进一步加强军队纪律，南昌行营决定在既有军纪约束之外，在前线部队中组织密查委员会，密查官兵有无不遵命令、营私舞弊、怠忽职守、勒索地方、招摇索贿、嫖赌吸毒酗酒等违纪行为。同时设立考验委员会，考核各级官兵作战和纪律情况，并依据考核成绩实施奖惩。\footnote{《剿匪军整顿军纪办法大纲》，《西路军公报》第9期，1934年3月15日。}蒋介石专电要求：“各团营连所派之采办不准其在地方民间自由购买物品，只准其在总指挥部所组织之采买组内采办。”\footnote{《蒋中正指示政治工作须深入匪区（1933年10月27日）》，蒋中正文物档案002010200096040。}力图杜绝部队乘机强买强卖。各部队也有相应的整饬军纪措施：“四十三师在宜黄设立粮食采办处以及提倡善良风俗移转社会风气等事，九四师党部办理官兵抚恤，九八师的救护队收容病兵纠察军纪抚恤难民，九九师由党部派员参加采买，使采买人不至压迫老百姓，和拿食盐来酬报抬伤兵的老百姓等。”\footnote{陈诚：《要认识领袖树立革命的中心力量》，《汗血月刊》第1卷第4期，1934年4月20日。}


根据对苏区民众“宽其既往，以广自新之路”\footnote{欧阳明：《匪区民众心理之分析与挽救》，《汗血月刊》第1卷第2、3期合刊，1934年2月20日。}的认识，1933年8月，南昌行营颁发《剿匪区内招抚投诚赤匪暂行办法》，规定对“投诚”的苏区一般人员可责成其父兄、邻右、房族长等具结担保领回，在家从事劳动，一年以内不准擅离所住区域。较重要“投诚”人员送感化院感化后再按上述程序处理。“投诚赤匪经该管县政府核准回原籍居住后，其生命财产应一律予以保护。”\footnote{《剿匪区内招抚投诚赤匪暂行办法》，《国民政府军事委员会委员长南昌行营处理剿匪省份政治工作报告》，第十三章，第5页。}1933年11月20日，南昌行营正式拟定《招抚投诚办法》公布。


其二，通过长期的“剿共”战争，国民党人逐渐意识到：“民心的向背，以利益为依归，我们要使民众归附我们，要使民众信仰我们的主义，空喊口号是没有用处，我们须从民众的实际利益加以维护。”\footnote{欧阳明：《匪区民众心理之分析与挽救》，《汗血月刊》第1卷第2、3期合刊，1934年2月20日。}为此，他们在经济上采取一系列措施，“复兴农村”，纾缓民生，动摇中共的民众基础。


土地问题是1930年代国共对立中一个十分吸引眼球的话题。面对中共土地问题上的积极政策，国民政府也亟思有所作为。蒋介石曾于1932年明确谈道：“对于乡村的土地问题，我们必须深刻留心才好，如果革命真正要成功的话，我们就是要平均地权，平均地权的实行，就是土地改革，中国所有一切问题，统统集中于土地问题上……要能切实做平均地权的工作，革命才有成功的胜算。”\footnote{蒋介石：《军队政治工作方法的改善》，《先总统蒋公思想言论总集》第10卷，第576页。}但是，对于平均地权的方法，他并不认同中共的土地分配做法，对农村土地实际占有状况的估计，也偏于乐观。1933年12月，蒋致电汪精卫，较为清楚地表达出其关于土地问题的立场：

\begin{quote}

今日中国之土地，不患缺乏，并不患地主把持，统计全国人口，与土地之分配，尚属地浮于人，不苦人不得地，惟苦地不整理。即人口繁殖之内地省区，亦绝少数百亩数千亩之地主，而三数十亩之中小耕农，确占半数以上。职是之故，中正对于土地政策，认为经营及整理问题，实更急于分配问题。既就分配而言，本党早有信条，即遵奉平均地权遗教，应达到耕者有其田之目的，而关于经营与整理，则应倡导集合耕作以谋农业之复兴。盖本党立场，不认阶级，反对斗争，关于土地分配，自应特辟和平途径，以渐进于耕者有田。\footnote{蒋介石：《电汪院长叶秘书长为解决土地问题详陈分配及整理意见以供参考》，《军政旬刊》第8期，1933年12月30日。}
\end{quote}


本着上述认识，国民政府一方面要求：“对于被匪分散之田地，有契据有经界者，以契据付审查，无契据有经界者，以证明书状付审查，办理完竣，一律发还原主”，\footnote{《鄂豫皖三省剿匪总司令部训令》，《西路军公报》第2期，1933年11月。}原则上承认和维护地主对土地的所有权；另方面，又要求地主对农民有一定的让步。1933年夏，南昌行营颁布《处理匪区土地、地租、田赋、债务办法》，规定凡1932年前的地租、田赋蠲免，债务缓还。同时实行《均耕法》，主要内容是：土地仍归原主，佃户受佃承耕，业主不得夺佃；凡有田两百亩以上者，课累进税，税金交农村复兴委员会支配；无主土地，由农村复兴委员会代管发佃，其地租亦由农村复兴委员会代收支配。根据这一法案内容，1933年10月，蒋介石亲电指示：“我军占领地方现在未收之谷子概归今年所种之佃户收获，以济贫农。明年再照土地条例妥为处理。地方人民从前所欠各种债务一律展期清理，其各债主不得追缴。”\footnote{《蒋中正电示周浑元处理占领地方谷物收获与地方人民所欠债务（1933年10月4日）》，蒋中正文物档案002010200095017。}1934年8月，再次电令新收复区当年之农产物，概归当年耕种者收获，原业主不得索取田租。10月，又下令收复区从前所欠田租、房租，均予免缴，其他债务延期清理，并规定应减免利息及最高利率限制。


作为具有强烈实用取向的领导人，蒋介石虽然主张渐进的土地整理，但占领苏区后具体处理土地问题时则采用了变通的办法。他明确指示：“如何处置土地，不一定要有呆板的方法，应当以补助剿匪进行为前提，因地制宜的去办，耕者有其田，平均地权，或者地还原主，或者实行二五减租，都是可以的；只要于剿匪进行有利，都可以斟酌办理。”\footnote{蒋介石：《清剿匪共与修明政治之道》，《先总统蒋公思想言论总集》第10卷，第623页。}因此，福建事变后，鉴于“前十九路军驻龙岩时，不分业佃，一律计口授田，现均有田可耕，确亦相安”，国民政府决定对“现在之承耕者计口授佃，不予变更”。\footnote{《国民政府军事委员会委员长南昌行营处理剿匪省份政治工作报告》，第九章，第40页。}等于默认了计口授田的现实。在原则维护地主土地所有权时，相当程度上考虑到普通农民的现实利益，对缓和农民不满情绪，抵消中共土地政策影响，不无意义。


针对“剿共”区域农村的破败状况，国民政府也采取了一些救济和保护措施。1933年4月，设立四省农民银行，蒋介石自兼董事长，展开金融救济农村活动。同时，本着“军民合作之方向，不得以军队便利为出发点，应以救济民众为出发点”\footnote{《剿匪军救济民众办法大纲》，《西路军公报》第9期，1934年3月15日。}之原则，蒋介石通令前方国民党军开展救济民众运动，要求展开以民生为基础的生命安全救济、生产救济、饥寒救济、教育救济等多项救济措施救济。1934年间南京政府拨发江西“剿匪”善后治本费200万元、治标费120万元，从治本费中提出30万元，加上农行商借的40万元共70万元用于救济工作，由江西省农村合作委员会负责办理。到1934年9月底，共放款358569.5元。\footnote{《江西省各收复县区办理农村救济组社员贷款一览表》，《国民政府军事委员会委员长南昌行营处理剿匪省份政治工作报告》，第九章，第32页。}为减轻民众税负，江西省政府决定实行“一税制”，即将各种税捐合成总数，冠以田赋税目统一收取，收取方式也由向就近钱粮柜所分期缴纳改为直接到县交纳，减少苛捐杂税和中间盘剥。这些措施或为治标，或系“杯水车薪”，\footnote{《国民政府军事委员会委员长南昌行营处理剿匪省份政治工作报告》，第五章，第43～44页。}但做与不做，其产生的政治影响，终究有别。


大力推广合作社，是蒋介石力图复兴农村经济的另一重要举措。合作社在20世纪二三十年代的中国颇具影响，一批社会力量投身于合作运动中，国共两党也都对之倾注热情。蒋介石强调：“农村合作事业，就是救济农村最紧要最要好的一个办法”；\footnote{《总理、总裁关于合作之言论》，《革命文献》第84辑，台北，中国国民党中央党史委员会，1980，第216页。}“发展农业，自以创设合作社为根本要图”。\footnote{《鄂豫皖三省剿匪总司令部训令》，《西路军公报》第2期，1933年11月。}1931年6月，国民政府实业部颁布《农村合作社暂行规程》，第一次以部令形式公布有关农村合作的章则。次年在湖北集训一批县级人员，为通过行政力量组织互助社提供干部。1933年10月和次年1月，南昌行营先后颁布《剿匪区内各省农村合作社条例》及《施行细则》、《剿匪区内农村合作委员会组织章程》等，要求“剿匪”区内各省设立农村合作委员会，大力推广合作运动。在江西将农村手工业合作改为“利用合作”，由政府贷给资金，各地利用当地的手工业特产，组织起来从事生产。到1934年1月，江西已成立496个合作社，有15000户社员，占当时全省总农业户400万户的0.375\%。\footnote{《江西省各县合作社种类数量及区联合会统计表》，《国民政府军事委员会委员长南昌行营处理剿匪省份政治工作报告》，第九章，第14页。}1934年底，进一步发展到1078个，\footnote{郑厚博：《中国合作社实况之检讨》，《实业部月刊》1936年10月号。}次年增加到2846个，社员231142户，股金1306369元，\footnote{孙兆乾：《江西调查日记》，萧铮主编《民国二十年代中国大陆土地问题资料》第171种，第85340页。}发展速度在全国名列前茅。行营同时提出，把拨给各县的善后经费，一律移充为当地农民加入合作社的股金，不得用于其他开支。国民党方面的报告自称：江西临川、崇仁、黎川等“六县中之设有合作社者，都有优良成绩”。\footnote{《检阅临川崇仁宜黄南丰南城黎川六县清乡善后事务之总讲评》，《军政旬刊》第26期，1934年6月30日。}


不过，和中共的合作社组织一样，国民党通过政权推动展开的合作社也具有较强的政治干预性质：“合作制度虽然是经济性质的，但却和政治性质的保甲制组织有直接的关系。联保主任或保长，常常是指派定的合作社的‘当然理事’。”\footnote{王志莘：《合作运动》，《中国经济年鉴（1936年）》，商务印书馆，1936，第881页。}合作社的经济效能发挥尚不充分。而政治推动由于经费、组织的限制又难以深入，所以，合作社发展虽然相对较快，但面对广大农民，其绝对值仍然很低，“仅能作微小之贡献”。成立起来的合作社也“恐不能得适合之管理”，\footnote{《国联专家视察江西建议书提要》，《军政旬刊》第23期，1934年5月31日。}难免为某些权势阶层控制、中饱。说起来，蒋介石当时对合作社其实还有更高的期待，他曾设想，通过合作社和农村复兴组织的推动，“各农村之田地，将陆续尽归农村利用合作社管理，而合作社全体社员，尽为农村田地之使用者，无复业主自耕农佃农雇农之分，则总理耕者有其田之主张，固不难具体实现，即彻底改良农业之方法，亦得以切实推行”。\footnote{《鄂豫皖三省剿匪总司令部训令（1932年11月）》，《西路军公报》第2期，1933年11月。}这样美好的愿望，却并没有切实可行的措施予以实现。


其三，进一步严密对苏区的全面封锁。中央红军前四次反“围剿”虽然均获胜利，但由此造成的资源损失也相当巨大，“围剿”后大批国民党部队继续环绕不退又限制着苏区的发展，军事的胜利不能完全掩盖苏区内部资源困乏及发展受限的危机。早在1931年中共方面其实就已意识到：“目前敌人尚未下绝大决心来封锁苏区，所以日常用品许多还可以入口。但我们要知道，阶级斗争日益尖锐和剧烈，敌人也就必然的更严密的来封锁苏区。我们为巩固政权，进攻敌人，在经济上须有充分的准备。”\footnote{《闽西苏维埃政府经济委员会扩大会议决议案（1931年4月15日）》，《革命根据地经济史料选编》（上），江西人民出版社，1986，第72页。}同理，1932年国民党方面也从中共被俘人员口供中推断：“匪区内除瑞金一县有少数货物买卖外，在他各县荒凉万分，若我中央能以此时一面给予政治上之宣传打击一面施坚壁清野封锁外物运入，则不出一年，不打自灭。”\footnote{李一之：《剿共随军日记》，第131页。}因此，国民党军有针对性地对苏区实施封锁，江西全省划分为8个封锁区，各区设监察官，由当地最高驻军长官担任，监察的层次分师、旅、团三级，均由各级军事长官兼任，政工人员担任巡查。半“匪区”、邻“匪区”各县，一律设立封锁匪区管理所，由县长兼任所长。水路方面颁布《整理赣江封锁计划大纲》，设立封锁赣江万（安）丰（城）间水道督察处及13个封锁管理所，加紧对赣江沿线的全面封锁。


1933年7月以后，南昌行营制定《粮食管理办法》、《合作社购销食盐办法》、《部队食盐采购办法》等一系列规章，规定粮食、食盐、火油、中西药品、布匹、服装、军用品、统销军用品、燃料等，以官督商办为原则，集中公卖，凭证购买，每口良民按定额供应。同时设立粮食管理处，食盐、火油管理处，交通管理处，负责组织实施对苏区物资的封锁。以县、区、保联主任及当地士绅组织各级公卖委员会，下设公卖处，负责购进、运销事宜。偷运或“济匪”者，轻者没收物资、罚款，重者判刑直至处死刑。


封锁制度严密化后，苏区物资供应相当紧张。尤其是维持正常生存必不可少的食盐被控制输入后，由于不能自产，极度紧缺。红军撤出中央苏区、国民党军进占瑞金后注意到：“居民食淡过久，肌肉黄瘦，殆不类人形。”\footnote{李渔叔：《瑞金匪祸记》，《收复瑞金纪事》，中国国民党陆军第十师特别党部，1935，第21页。}其残酷情形可见一斑。米夫在共产国际的发言中谈道：“过去苏区与国统区的经济往来是相当容易的，而现在则要困难几十倍。国民党对于跟苏区经商的商人不惜采取各种镇压、枪毙的手段，而这种封锁产生了效果。如果说过去不管怎样总能把一些工业品带到苏区来，而现在可能性极少。”\footnote{《共产国际执行委员会政治书记处会议速记记录（1932年12月11日）》，《共产国际、联共（布）与中国革命档案资料丛书》第13册，第274页。}与此同时，临近苏区的国民党控制区民众也深受封锁之苦：“食盐公卖以后，各地时有被少数甲长操纵渔利之事，规定每人只购四两，但与甲长关系密切者，至少可买四斤，与甲长无私情者，即四两亦不能到手。”\footnote{《为各保甲长操纵公卖仰转饬严究》，《军政旬刊》第1期，1933年11月20日。}


其四，严密政治组织，加强行政控制。国民党从组织上看，不是一个十分严密的政党，其地方自治的治理原则也从理论上限制着政权垂直权力的过度伸展。而蒋介石以军事强人控制政权的现实及南京政府对全国实际统治力的薄弱又进一步影响着南京政权行政控制力的发挥。在与中共对垒过程中，蒋介石得出结论：中共严密的组织使其具有强大的群众动员能力，是中共能将自身实力充分发挥的一个关键原因，而国民党政权自身则“政府自政府，人民自人民，军队自军队，各不相谋，甚至省政府和县政府之间，也不能十分联络得好”。\footnote{蒋介石：《剿匪基本工作之研究》，《先总统蒋公思想言论总集》第11卷，第38页。}针对此，蒋力图在现有框架内对政治组织有所改进，加强政权的行政和社会控制力。


按照南京政府的地方政治体制，省直管县，传统中国原有的州府一级机构被取消，这样的政治结构源于孙中山倡导的地方自治原则。但是，当时省境庞大，交通不便，省对县的管理常常鞭长莫及。赣南的三南地区（全南、定南、龙南三县）距省会南昌有六七百公里之遥，又无公路可通，省级管理几乎无从措手。1932年8月，鉴于现有省县机构难以满足管理需要的实际状况，国民政府行政院和豫鄂皖三省“剿匪”总司令部决定在省和县之间按区域增设派出性的专员公署，加强省、县之间联系。1933年1月，南昌行营改组江西各地专员公署，扩大专员公署权力，专员一律兼任该区保安司令，并须逐渐兼任专员公署所在地县长；同时，为专员公署增加经费，配设技术人才。1934年福建事变后，又在福建推行专员公署制度，将全省分为十区，设置专员。专员公署的建立健全，在省、县之间增加了一级具有相当行政权的督察机关，可以考核地方官员、查核地方财政、核转省县间往来文件。每个专员公署一般管辖6个县，这对交通不便状况下行政权力的有效施展不无裨益。


除设立专员公署外，南昌行营对一些离县治较远，几县交界且“平时政治力量，已有鞭长莫及之患”\footnote{《筹设找桥慈化洋溪三特别区政治局案》，《军政旬刊》第11期，1934年1月30日。}地区，本着“适应剿匪需要，增加行政效率”\footnote{《国民政府军事委员会委员长南昌行营办法江西特别区政治局组织条例》，《西路军公报》第11期，1934年5月15日。}的原则，加设特别行政区，就近控制。1933年7月27日，南昌行营在江西藤田等四处设置特别行政区，以永丰、乐安、吉水三县交界的一部分治地为藤田特别行政区，祟仁、宜黄两县交界的一部分治地为凤冈特别行政区，同时设立新丰、龙岗特别行政区（旋被裁撤）。10月、12月，行营又分别在宜丰找桥、宜春慈化和安福洋溪、井冈山大汾设立特别行政区。特别行政区设政治局，隶属于行营和江西省政府，负责“处理全区一切行政事务”；“政治局对各级机关之关系，与县政府同”。\footnote{《江西省特别区政治局组织条例》，《国民政府军事委员会委员长南昌行营处理剿匪省份政治工作报告》，第一章，第28页。}


乡村保甲制度是中国农村实施已久的传统控管体系，民国成立后一度被废弃。1930年代初，为适应“剿共”军事的需要，加强对基层社会的控制，南京国民政府开始重建乡村保甲。1932年8月，鄂豫皖三省“剿匪”总司令部颁布《剿匪区内各县编查保甲户口条例》，开始大面积推行保甲制度。1933年9月，江西81县中有62县进行了保甲编组，保甲网络基本成型。蒋介石对保甲十分重视，指示：“进剿时亲查保甲，帮修寨碉，最为重要。并择其要冲之地遴选稳实保甲长，给其能足自卫之旧枪，派能干可靠之官兵数人训练监督……使其自守，则为根本之治也。”\footnote{《蒋中正电示樊崧甫进剿时亲查保甲帮修寨碉慎选保甲长给枪自卫为要（1933年5月29日）》，蒋中正文物档案002010200084061。}“进剿”期间，保甲的建设和恢复在数据上继续显示出较快的进展，据江西省民政厅统计，到1935年江西各县编成26584保，259066甲。保甲制度的推展，对国民党权力体系的垂直延伸及反制中共的发展、渗透有着一定作用。

\subsection{“七分政治”的效用}


国民党政权在政治、经济、社会上与中共全面对垒的一系列对策，取得了一定效果。国民党军军风纪在此期间确有改善，国民党方面各级人员不约而同谈道：“从前没有组织运输队的时候，民众因为怕拉夫的关系，军队所到的地方，逃避一空，现在不同了，我们军队所过的地方，老百姓排班的站着，送茶送水的，络绎于途，请向导也不发生困难，这是由于军纪严明的效果。”\footnote{欧阳明：《匪区民众心理之分析与挽救》，《汗血月刊》第1卷第2、3期合刊，1934年2月20日。}诸如强奸、抢劫等恶性案件明显减少。当时，苏区有“有一部分群众脱离政府，如东陂、黄陂、吴村有一部分群众听反动派造谣，见红军到逃跑上山，反接济白军靖匪的食粮”\footnote{《江西省苏维埃政府训令（第九号）（1933年4月5日）》，《江西革命历史文件汇集（1933～1934年）》，第87页。}；一些地区“群众大受敌人欺骗，反水成立守望队替敌人担任秘密通讯员等”。\footnote{《中共湘赣省委工作报告》，《湘赣革命根据地》（上），第491～492页。}对此，蒋介石曾不无得意地宣称：

\begin{quote}

从前我们军队到的时候，一般人民一定都被土匪裹挟去；现在我们这一次到了棠阴，一般人民却希望土匪早一些快走，我们军队一到，他们就出来，可见赤匪的手段无论怎么毒辣巧妙，无论对于部下对于人民监督怎么严密，在最短期间或可稍微发生效力，但用了一年半载以后，一概无用，而且还要发生反结果。\footnote{蒋介石：《军人精神教育之要义》（一），《庐山训练集》，第171页。}
\end{quote}


不过，国民党方面采取的这一系列政策，似也不应予以过高估计，从总体上看，所谓“七分政治”其实远远未能达到超越于军事之上的效果，在国民党政治理念、统治方式、统治基础不可能根本改变的背景下，其推出的许多措施实际效果有限。当时江西金溪县长朱琛上书蒋介石谈道：“处处离开民众，任何良法适得其反。故保甲造成土劣集团，保卫团成为地痞渊薮，建筑堡垒，徒劳民财，演成政府求治之心益切，而人民所受之痛苦则日深。其原因均为政烦赋重，处处予贪污土劣剥削之机会，故欲扬汤止沸，莫如釜底抽薪也。”\footnote{《朱琛上蒋介石意见书》，傅莘耕：《金溪匪区实习调查报告》，萧铮主编《民国二十年代中国大陆土地问题资料》第173种，第86111页。}这一议论确非虚言。当国民党军实行普遍的动员贯彻其全面战争计划时，对民众的压力也在加大。


第五次“围剿”中，由于碉堡、公路的大规模修筑及军队的增加，民众的劳役和供给压力较之前几次继续增加。当时有报告谈到修筑碉堡的费用：“只以取沙而论，其工料每堡需千六七百元，每碉需七八百元。假定赣属沿江构筑百座，堡占三分之一，碉占三分之二，预算不下十万元。”\footnote{《某人电蒋中正据张达称构筑赣江及湘粤赣边境各碉堡预算甚巨》，蒋中正文物档案002090300073351。}“金溪县建筑堡垒五，碉堡八，共费万元以上。”\footnote{傅莘耕：《金溪匪区实习调查报告》，萧铮主编《民国二十年代中国大陆土地问题资料》第173种，第86128页。}更成问题的是，军队和官员借修碉堡盘剥民众，进一步加重民众负担：“各部队修筑碉堡与工事，各种器具，多系借自民间，损坏既无赔偿，移防辄多携去。且士兵亲往民家搜借，更难保无违反纪律行为，影响军民恶感甚大。”\footnote{《政训处长贺衷寒等转呈金文质呈拟对民众工作意见》，《军政旬刊》第13、14期合刊，1934年2月28日。}有些部队“甚至将避逃堡内之民众，悉数赶出堡外，以致流为匪用”。\footnote{《北路总司令部赣浙闽皖边区警司令部转报余江县一带党政军情形》，《军政旬刊》第7期，1933年12月20日。}而一些地方官员则把构筑碉堡当作聚敛手段，江西萍乡北一区区长为建筑碉堡，“于地方筹集洋二万余元”，挪用寺庙“砖木值洋约五百元并不给价”。\footnote{《萍乡北一区福兴寺僧蔡鑫等具诉区长萧造时假公济私强挑建碉》，《西路军公报》第7期，1934年1月15日。}莲花县“建筑公路，各区各保，咸以摊派方式担任，不第工资无着，且须自备伙食……此外如派筑碉堡（闻派筑碉堡亦不供食不给资且砖石亦系按人摊派送去）架设电话，与服役于军队之运输，一般壮年男女劳役几无暇日”。\footnote{赵可师：《赣西收复区各县考察记》，《江西教育旬刊》第10卷第4、5期合刊，1934年7月11日。}


从收揽民心出发，第五次“围剿”期间，蒋对土豪劣绅不无抑制之意，但由于国民党政权基层控制力薄弱，乡村治理重心很大程度上仍依赖士绅，蒋介石本人对乡绅和土地问题处理又首鼠两端，既不想失去士绅的支持，又想讨好民众，他在日记中说得很明白：“对于土地问题二说，一在恢复原状，归还地主；一在设施新法，实行耕者有其地主义。对于耆绅亦有二说，一在利用耆绅，招徕士民；一在注重贫民，轻视耆绅，以博贫民欢心。余意二者可兼用也。”\footnote{《蒋介石日记》，1932年6月7日。}在此背景下，地方官员对乡绅常常是投鼠忌器，涉及乡绅案件，“地方县政，往往迁延顾忌，久不处理，以致民众感受痛苦，来部指控者，纷至沓来”。\footnote{《李蕴珩致蒋介石鱼电》，《军政旬刊》第3期，1933年11月10日。}而在收复地区土地处理这样关乎利益取向的重大问题上，蒋介石也经常倾向乡绅，认为该事项“端赖地方士绅相助”，并制定《敦促各县士绅回县参加清乡善后工作办法》，要求各地方政府“劝导各地士绅从速回乡，共襄要政”。\footnote{《国民政府军事委员会委员长南昌行营处理剿匪省份政治工作报告》，第九章，第39～40页。}对此，曾有人明确指出：“各县乡村受共产党‘有土皆豪，无绅不劣’之宣传，民众对一般绅士，已无信仰，或且憎恶，现又加以劣绅回乡工作，益增人民仇恨，殊失政府招回公正士绅之德意，而公正士绅益将裹足不前，乡村政治更不堪闻问矣。”\footnote{《新丰特别区政治局局长刘千俊报告匪区民众根本动摇情形匪之维持残局原因及所拟对策》，《军政旬刊》第7期，1933年12月20日。}由于乡绅本为强势一方，两头讨好的结果事实上仍然是对乡绅的纵容。


第五次“围剿”中，国民党军军纪虽有改善，但违犯军纪、侵害百姓的事件仍时有发生。1933年8月，到江西调查的地政调查人员遇到民众诉说：“第六师士兵强行贱买谷米，只付一元要挪一石，欲阻不得，反被恶骂毒殴。”\footnote{傅莘耕：《金溪匪区实习调查报告》，萧铮主编《民国二十年代中国大陆土地问题资料》第173种，第86130页。}12月，国民党军第十二师通令全师指出：“近日以来，纪律废弛，日甚一日，如各部采买，藉口纸币不能换散，强压购买，或以地方匪化心理，擅取民物，不给代价，更有兵到之地，翻箱倒箧，形同洗劫！以致民怨沸腾，军誉扫地。”\footnote{《令各部队严整军纪（1933年12月21日）》，《汗血月刊》第1卷第2、3期合刊，1934年2月20日。}强买强卖仍是痼疾。军队拉夫也不能完全禁绝，南昌行营别动队总队长康泽报告，江西余江一带驻军“每于开拔时，无论民夫递步哨，不分皂白，一律拉用，甚至哨丁携带公文，中途相遇，均遭强迫拉去，并有凶殴情事”。\footnote{《北路总司令部赣浙闽皖边区警司令部转报余江县一带党政军情形》，《军政旬刊》第7期，1933年12月20日。文中“民夫递步哨”应为“民夫与步哨”之解。}占用民房更是司空见惯：“士兵所驻之房屋，多半是人民住定之内室，责令空出一部分。其厅堂及厨房既归士兵占有，房东不得自由。彼去此来，继续无间，人民住宅几成兵站。至于借用日常需要器具，任意损坏或遗失，毫不负责，犹其余事。”\footnote{《训令本部军队所到境地非公地不敷住宿时勿得轻驻民房》，《西路军公报》第10期，1934年4月15日。}直到“围剿”末期，扰民现象始终不绝，陈诚在家书中写道：“各部军风纪亦不甚好，如第四、第八十八师到处杀牛杀鸡、挖蕃薯花生等等，人民被扰情形亦可想见。”\footnote{《至各处视察沿途死病士兵夫不知其数各部军风纪亦不甚好（1934年10月13日）》，《陈诚先生书信集·家书》（上），第283页。}第十师攻占瑞金当天，就有官兵拒不听令，“违令下乡，而且短给物价”。\footnote{孙一鲲：《匪都“瑞京”破灭的因果》，《收复瑞金纪事》，陆军第十师特别党部，1935，第61页。}军队建设的种种问题，蒋介石历历在目，日记中有清楚展现：“上午由南昌来抚州，沿途见军队之污秽，与人民之痛苦，伤心自罪，不知所止。”\footnote{《蒋介石日记》，1933年11月28日。}


至于政权建设，当时调查报告承认：“赣省年来提倡廉洁政治，绝少贪污，风气确已丕变，惟此亦不过限于县长而已，至于各区区长，仍多属贪污土劣之大集团。”\footnote{《检阅临川崇仁宜黄南丰南城黎川六县清乡善后事务之总讲评》，《军政旬刊》第26期，1934年6月30日。}亲身指挥战事的陈诚也坦白指出：“此次集合数省兵力，大举围剿，并于地方党政设施，妥为规划，期以政治力量，摧毁赤匪根基。深谋远虑，纤细靡遗，匪患削平，指顾可待。惟查各县区地方，对于所颁法令，未能切实遵行，即以编查保甲团队，封锁匪区，及构筑碉堡机场公路诸端而论，或旷日废时，一无所就；或有名无实，粗具规模；便利剿匪之实效难期，徒苦人民之弊病已见。”\footnote{《函呈动员民众参加剿匪（1933年10月24日）》，《陈诚先生书信集·与蒋中正先生往来函电》（上），第116页。}平心而论，政治的改变本非一朝一夕所能达成，何况战争环境下各种机制均待健全，政治能够不拖军事后腿就已可算成功，出现种种问题实无足怪。只是，这一现实显然难符蒋介石的期望，所以他在1934年底曾感慨叹息：“从前我们讲剿匪要有三分军事，七分政治，这不过是一种理想，事实上直到现在不过三分政治，仍旧七分军事！最多也只是军事政治各半！”\footnote{蒋介石：《剿匪胜利中吾人应继续努力》，《先总统蒋公思想言论总集》第12卷，第593页。}所谓“三分军事，七分政治”，并不一定是实际效果的真实反映，更多体现的还是蒋对政治“剿共”的重视。

